\chapter{Conclusão} \label{cap:Conclusao}

Este trabalho abordou o desafio de otimizar o balanceamento de carga em Redes de Distribuição de Conteúdo (CDNs) para \textit{streaming} de vídeo adaptativo, propondo um algoritmo baseado em monitoramento de memória principal e taxa de leitura de disco dos servidores. A motivação central foi superar as limitações de estratégias convencionais como \textit{Round-Robin} e Seleção Aleatória, que operam sem considerar o estado real dos recursos dos servidores, resultando em distribuição ineficiente de carga e degradação da Qualidade de Experiência (QoE) dos usuários, especialmente em ambientes com servidores de distribuição de conteúdo com capacidades heterogêneas.

Através de implementação completa e validação experimental sistemática em quatro cenários distintos, variando escala (100 e 1000 usuários) e heterogeneidade dos servidores (homogênea e heterogênea), este trabalho demonstrou que o monitoramento inteligente de recursos permite melhorias substanciais em QoE. Este capítulo sintetiza as principais descobertas, valida as hipóteses estabelecidas, discute as contribuições do trabalho, reconhece suas limitações e aponta direções para pesquisas futuras.

\section{Síntese dos Resultados}

Os experimentos revelaram melhorias dramáticas quando o algoritmo proposto foi aplicado em cenários com servidores de capacidades heterogêneas sob carga significativa. No Cenário 1 (100 usuários, servidores heterogêneos), o algoritmo com intervalo de atualização das informações de uso da memória volátil e de taxa de leitura da memória secundária (disco) de 6 segundos reduziu a latência média em 85,7\% comparado ao \textit{Round-Robin}, de 2345,5ms para 336,6ms, enquanto manteve qualidade de vídeo consistentemente em 1080P e eliminou travamentos.

No Cenário 3 (1000 usuários, servidores heterogêneos), o mais desafiador dos experimentos, os resultados foram ainda mais contundentes. Enquanto estratégias convencionais colapsaram completamente — com \textit{Round-Robin} registrando 7 travamentos (15,12 segundos totais) e Seleção Aleatória registrando 8 travamentos (29,72 segundos totais), ambos com qualidade degradada para \textit{144P} — o algoritmo proposto com intervalo de 6 segundos eliminou totalmente os travamentos, manteve qualidade predominantemente em \textit{720P}-\textit{1080P}, e reduziu a latência em 70,0\% (de 3988,4ms para 1197,2ms).

O algoritmo de balanceamento de carga baseado no monitoramento de memória, mesmo com o intervalo de atualização de 30 segundos, representando dados moderadamente desatualizados, demonstrou robustez notável. No Cenário 3, apresentou apenas 4 travamentos curtos (total de 6,42 segundos, média de 1,6s cada), ainda substancialmente superior às estratégias convencionais, validando a capacidade do mecanismo de interpretação de dados desatualizados.

Em contraste, nos cenários com servidores homogêneos (Cenários 2 e 4), todas as estratégias convergiram para desempenho similar, confirmando que o valor do monitoramento inteligente é proporcional à heterogeneidade do ambiente. Esta descoberta define claramente os contextos onde o algoritmo proposto oferece benefícios tangíveis versus contextos onde estratégias simples são adequadas.

\section{Validação das Hipóteses}

As hipóteses estabelecidas no Capítulo~\ref{cap:Proposta} foram integralmente validadas pelos resultados experimentais:

\textbf{Hipótese 1:} \textit{O algoritmo de balanceamento baseado em monitoramento de memória apresentaria desempenho superior aos métodos \textit{Round-Robin} e Aleatório em cenários heterogêneos.}

\textbf{Validação:} Confirmada inequivocamente. Reduções de latência de 70-85\%, eliminação ou drástica redução de travamentos (de 7-8 para 0-4), e manutenção de qualidade alta e consistente (720P-1080P vs. 144P) demonstram superioridade clara nos Cenários 1 e 3.

\textbf{Hipótese 2:} \textit{O intervalo de atualização de 6 segundos forneceria melhores resultados que o intervalo de 30 segundos.}

\textbf{Validação:} Confirmada, com ressalva importante. O intervalo de 6 segundos efetivamente apresentou o melhor desempenho absoluto em todos os cenários heterogêneos. No entanto, o resultado mais significativo foi que o intervalo de 30 segundos, apesar de usar dados cinco vezes mais antigos, ainda superou substancialmente as estratégias convencionais, demonstrando robustez excepcional do mecanismo de interpretação de dados de uso de memória volátil e taxa de leitura de memória secundária moderadamente desatualizados.

\textbf{Hipótese 3:} \textit{O algoritmo seria robusto mesmo com informações moderadamente desatualizadas, operando eficientemente com intervalos de atualização maiores.}

\textbf{Validação:} Confirmada de forma conclusiva. No Cenário 1, o intervalo de 30 segundos manteve redução de 74,2\% na latência sem travamentos. No Cenário 3, embora tenha apresentado 4 travamentos, seu desempenho permaneceu superior à Seleção Aleatória (8 travamentos, 29,72s) e comparável ao \textit{Round-Robin} em número de eventos, mas com duração significativamente menor (6,42s vs. 15,12s).

Estas validações demonstram que o algoritmo não apenas funciona conforme projetado, mas excede expectativas em termos de robustez, oferecendo \textit{graceful degradation} quando opera com informações desatualizadas, característica essencial para escalabilidade em sistemas de larga escala.

\section{Robustez do Algoritmo Proposto}

Um resultado particularmente significativo e promissor para aplicações práticas é a robustez do algoritmo quando operando com dados relativamente desatualizados. O desempenho mantido com intervalo de atualização de 30 segundos valida o conceito de \textit{Load Interpretation} adaptado das equações propostas, onde as chegadas esperadas durante o período de desatualização são incorporadas no cálculo probabilístico de seleção de servidores.

Esta robustez tem implicações práticas cruciais para escalabilidade. Sistemas de larga escala com centenas ou milhares de servidores enfrentam custos significativos de monitoramento: sobrecarga de processamento para coleta de métricas, largura de banda para comunicação via MQTT, latência de propagação de informações, e complexidade operacional. A capacidade de operar eficientemente com intervalos de atualização maiores reduz todos esses custos proporcionalmente.

Adicionalmente, a robustez a dados desatualizados proporciona tolerância natural a falhas transitórias no sistema de monitoramento. Se o \textit{broker} MQTT sofrer latência temporária ou um servidor não publicar métricas momentaneamente, o balanceador continua operando com informações ligeiramente mais antigas, mas ainda eficazmente, evitando degradação catastrófica.

\section{Fatores Determinantes do Desempenho}

A análise dos resultados revela dois fatores principais que determinam quando e quanto o algoritmo proposto proporciona valor: heterogeneidade dos servidores e nível de concorrência.

\subsection{Heterogeneidade como Fator Crítico}

Os resultados demonstram inequivocamente que a heterogeneidade dos servidores é o fator determinante primário para o valor do balanceamento inteligente. A divergência dramática de desempenho entre cenários heterogêneos (Cenários 1 e 3) e homogêneos (Cenários 2 e 4) não deixa dúvidas: monitoramento e adaptação são essenciais quando servidores possuem capacidades diferentes, mas oferecem valor marginal quando todos são equivalentes.

Esta conclusão tem implicações importantes para a indústria de CDNs e \textit{edge computing}, onde heterogeneidade de disponibilidade de recursos é a norma, não a exceção. Embora ambientes virtualizados modernos frequentemente utilizem servidores com capacidades físicas homogêneas, a heterogeneidade manifesta-se através da variação dinâmica de carga de trabalho. Servidores em um cluster não se dedicam exclusivamente a servir conteúdo de vídeo — simultaneamente executam outras aplicações, processam tarefas de \textit{background}, respondem a requisições de diferentes \textit{tenants} em ambientes multi-inquilino, e lidam com processos do sistema operacional. Esta concorrência por recursos cria heterogeneidade dinâmica: mesmo servidores idênticos em \textit{hardware} apresentam disponibilidade drasticamente diferente de memória, \textit{CPU}, e largura de banda de disco a cada instante, dependendo das demandas concorrentes que enfrentam.

Nestes contextos, que representam a realidade operacional da maioria dos sistemas de distribuição de conteúdo em produção, estratégias simples como \textit{Round-Robin} são fundamentalmente inadequadas. Ao distribuir uniformemente a carga, efetivamente aplicam a mesma demanda a servidores com capacidades drasticamente diferentes, resultando em servidores com menor capacidade saturados enquanto servidores com maior capacidade permanecem ociosos. O algoritmo proposto, ao observar e reagir à carga real, naturalmente direciona mais requisições para servidores capazes e menos para servidores limitados, aproximando-se de uma distribuição ótima.

\subsection{Concorrência como Teste de Estresse}

O nível de concorrência funcionou como teste de estresse efetivo, revelando as limitações das estratégias convencionais. Sob carga moderada (100 usuários), mesmo estratégias subótimas conseguem funcionar razoavelmente porque há capacidade suficiente no sistema para absorver ineficiências. Sob alta carga (1000 usuários), a ineficiência na utilização de recursos manifesta-se dramaticamente na forma de colapso de qualidade e emergência de travamentos.

O fato de que o algoritmo proposto manteve desempenho superior mesmo sob estresse extremo (Cenário 3) demonstra não apenas eficácia em condições normais, mas também resiliência sob picos de demanda. Esta característica é fundamental para sistemas de produção, que devem lidar graciosamente com variabilidade de carga típica de eventos ao vivo, lançamentos de conteúdo popular, ou padrões sazonais de acesso.

A análise conjunta destes fatores sugere uma regra prática: o valor do monitoramento inteligente é proporcional ao produto da heterogeneidade pela escala. Em ambientes homogêneos sob baixa carga, estratégias simples são adequadas. Conforme heterogeneidade ou concorrência aumentam, os benefícios do monitoramento crescem. Quando ambos são altos, como no Cenário 3, o monitoramento torna-se não apenas benéfico, mas essencial para operação funcional.

\section{Contribuições do Trabalho}

Este trabalho oferece múltiplas contribuições para o estado da arte em balanceamento de carga para CDNs e \textit{streaming} de vídeo adaptativo:

\subsection{Algoritmo de Balanceamento Baseado em Memória}

Proposta e implementação de um algoritmo de balanceamento de carga que utiliza consumo de memória principal e taxa de leitura de disco como métricas primárias para decisões de encaminhamento para atendimento às requisições de conteúdo. Diferentemente de abordagens convencionais que monitoram CPU ou conexões ativas, este trabalho sugere que métricas de memória e leitura de disco podem ser mais apropriados para servidores MPEG-DASH, refletindo diretamente a capacidade de servir segmentos de vídeo eficientemente.

\subsection{Mecanismo de Interpretação de Dados Desatualizados} 

Desenvolvimento de um mecanismo de \textit{Load Interpretation} que permite ao balanceador operar eficientemente mesmo com informações moderadamente desatualizadas. Através do cálculo das chegadas esperadas durante o período de desatualização, o algoritmo ajusta probabilidades de seleção para compensar a defasagem temporal, proporcionando robustez essencial para escalabilidade em sistemas de larga escala.

\subsection{Validação Experimental Abrangente}

Condução de experimentação sistemática e abrangente, com quatro cenários distintos testando as combinações de concorrência (baixa/alta) e heterogeneidade (homogênea/heterogênea), validando o algoritmo em condições diversas e identificando precisamente os contextos onde oferece maior valor. A metodologia experimental, com coleta de múltiplas métricas de QoE e análise quantitativa detalhada, estabelece base sólida para as conclusões.

\subsection{Demonstração da Heterogeneidade como Fator Crítico}

Identificação e demonstração clara de que a heterogeneidade dos servidores é o fator determinante primário para eficácia de balanceamento inteligente. Esta descoberta tem implicações práticas diretas para decisões de implantação e arquitetura de CDNs, orientando quando investir em monitoramento versus quando estratégias simples são suficientes.

\subsection{Implementação de Código Aberto}

Desenvolvimento de uma implementação completa e funcional utilizando tecnologias modernas (\textit{C\#} \textit{ASP.NET Core}, \textit{Rust}, \textit{Python}, \textit{Docker}, MQTT), com código organizado, documentado e disponibilizado de forma aberta. Esta contribuição permite reprodução dos experimentos, validação independente dos resultados, e serve como base para trabalhos futuros e adaptações práticas.

\subsection{Infraestrutura Experimental Reutilizável}

Criação de uma infraestrutura experimental completa incluindo servidores MPEG-DASH configuráveis, cliente de validação com visualização em tempo real, gerador de carga parametrizável, painel de monitoramento, e scripts de geração de gráficos. Esta infraestrutura pode ser reutilizada por pesquisadores e profissionais para validar outros algoritmos ou estudar diferentes aspectos de sistemas de \textit{streaming}.

\section{Limitações}

Embora os resultados sejam fortemente positivos e as contribuições significativas, reconhecem-se limitações importantes que devem ser consideradas ao interpretar os achados e planejar trabalhos futuros.

\subsection{Ambiente Experimental Virtualizado}

Os experimentos foram conduzidos em ambiente virtualizado utilizando \textit{Docker} em máquina única. Embora esta abordagem permita controle preciso sobre condições experimentais e reprodutibilidade, não captura completamente efeitos de rede real como variação de latência geográfica, perda de pacotes, jitter, ou congestão de rede. Servidores distribuídos geograficamente experimentam condições de rede heterogêneas que podem interagir com decisões de balanceamento de formas não observadas neste estudo.

\subsection{Escopo de Cenários de Concorrência}

Os experimentos focaram em dois níveis específicos de concorrência: 100 e 1000 usuários simultâneos. Embora esta escolha permita comparação clara entre baixa e alta carga, sistemas de produção experimentam gama muito mais ampla de cargas, desde dezenas até milhões de usuários. Adicionalmente, padrões de acesso real incluem crescimento gradual, picos abruptos, oscilações cíclicas (padrões diurnos), e cargas flash (eventos virais), comportamentos não testados sistematicamente.

\subsection{Algoritmo ABR Único}

A validação utilizou exclusivamente o algoritmo de seleção adaptativa de qualidade (\textit{Adaptive Bitrate - ABR}) padrão da biblioteca \texttt{dash.js}. Diferentes algoritmos ABR podem reagir de formas distintas às variações de taxa de transferência resultantes de diferentes estratégias de balanceamento, potencialmente alterando os resultados de QoE. A interação entre balanceamento de carga e algoritmos ABR é área complexa que merece investigação mais profunda.

\subsection{Métricas de Monitoramento Limitadas}

O algoritmo atual utiliza apenas duas métricas: consumo de memória principal e taxa de leitura de disco. Embora estas tenham se mostrado efetivas e apropriadas para servidores MPEG-DASH, sistemas reais enfrentam múltiplos tipos de gargalos. Utilização de \textit{CPU}, largura de banda de rede, latência de resposta, e localização geográfica são fatores adicionais que poderiam melhorar decisões de balanceamento. A otimização multi-objetivo com múltiplas métricas permanece inexplorada.

\subsection{Conteúdo de Vídeo Único}

Os experimentos utilizaram um único arquivo de vídeo codificado em múltiplas qualidades. Vídeos com diferentes características (ação intensa vs. cenas estáticas, resolução \textit{4K}, \textit{codecs} \textit{H.265}/\textit{AV1}, etc.) podem apresentar padrões de acesso e requisitos de recursos distintos, potencialmente impactando eficácia relativa das estratégias de balanceamento.

\subsection{Disponibilidade Garantida de Conteúdo}

Este trabalho assume que todos os servidores possuem cópia completa do conteúdo servido, simplificando o problema de balanceamento para apenas decidir qual servidor utilizar para cada requisição. Em CDNs reais, especialmente aquelas com grande catálogo de conteúdo, nem todos os servidores mantêm cópias de todo o conteúdo disponível. Estratégias de \textit{cache} parcial, replicação seletiva baseada em popularidade, e geo-distribuição de conteúdo introduzem complexidade adicional: o balanceador deve considerar não apenas a capacidade de recursos do servidor, mas também se o servidor possui o conteúdo requisitado em \textit{cache}.

Cenários onde conteúdo não está garantidamente disponível em todos os servidores requerem coordenação entre decisões de balanceamento e políticas de \textit{cache}, potencialmente introduzindo trade-offs entre escolher servidor com melhor capacidade disponível (que pode precisar buscar conteúdo no servidor origem) versus servidor com menor capacidade mas que já possui conteúdo em \textit{cache}. Esta limitação representa oportunidade significativa para extensão do trabalho, explorando balanceamento \textit{content-aware} que considera conjuntamente estado de recursos e disponibilidade de \textit{cache}.

\section{Trabalhos Futuros}

Os resultados deste trabalho e as limitações identificadas abrem múltiplas direções promissoras para pesquisas futuras.

\subsection{Implantação em Ambientes Distribuídos Reais}

Expansão dos experimentos para ambientes distribuídos geograficamente, com servidores em múltiplos \textit{data centers} e regiões. Isto permitiria avaliar como latências de rede reais, variações de largura de banda, e perdas de pacotes interagem com decisões de balanceamento, e como incorporar métricas de rede (latência, perda, jitter) nas decisões de roteamento.

\subsection{Integração com Sistemas CDN de Produção}

Validação do algoritmo em CDNs de produção servindo tráfego real de usuários, permitindo avaliar desempenho sob padrões de acesso autênticos, diversidade de dispositivos e condições de rede, e requisitos operacionais de sistemas em larga escala. Parcerias com provedores de CDN ou plataformas de \textit{streaming} seriam valiosas para esta validação.

\subsection{Aprendizado de Máquina para Predição de Carga}

Incorporação de técnicas de \textit{machine learning} para predição de carga futura dos servidores, permitindo decisões proativas ao invés de apenas reativas. Modelos de séries temporais (\textit{LSTM}, \textit{Prophet}) poderiam prever tendências de carga com minutos de antecedência, permitindo balanceamento preditivo que antecipa sobrecargas antes que ocorram.

\subsection{Otimização Multi-Métrica}

Expansão do algoritmo para considerar múltiplas métricas simultaneamente: memória, disco, \textit{CPU}, largura de banda de rede, latência geográfica para o usuário. Desenvolvimento de função objetivo multi-critério que pondere estes fatores apropriadamente, potencialmente utilizando técnicas de otimização multi-objetivo (\textit{Pareto}) ou aprendizado por reforço para aprender pesos ótimos dinamicamente.

\subsection{Balanceamento Cross-data center}

Extensão do algoritmo para balanceamento não apenas entre servidores de um único \textit{data center}, mas entre múltiplos \textit{data centers} geograficamente distribuídos. Isto introduz trade-offs complexos entre capacidade de recursos (servidor menos carregado pode estar longe) e proximidade geográfica (servidor próximo pode estar mais carregado), requerendo modelagem sofisticada de custos multi-dimensionais.

\subsection{Adaptação para Edge Computing e 5G}

Exploração de como o algoritmo pode ser adaptado para ambientes de \textit{edge computing} e redes \textit{5G}, onde recursos computacionais estão distribuídos em torres de celular e pontos de presença de borda. Nestes contextos, mobilidade dos usuários, \textit{handoffs} entre células, e latências ultra-baixas introduzem desafios adicionais.

\subsection{Estudo de Interação com Algoritmos ABR}

Investigação sistemática de como diferentes algoritmos \textit{ABR} (\textit{BOLA}, \textit{MPC}, \textit{Pensieve}) interagem com estratégias de balanceamento de carga. Exploração de \textit{co-design} onde balanceador e \textit{ABR} cooperam ativamente, com balanceador informando capacidade esperada e \textit{ABR} ajustando decisões de qualidade proativamente.

\subsection{Implementação de \textit{Cache} Inteligente Integrado}

Integração de estratégias de \textit{cache} inteligente que considerem não apenas popularidade de conteúdo mas também capacidade de recursos dos servidores. Servidores mais capazes poderiam \textit{cache}ar conteúdo de maior demanda, enquanto servidores limitados focariam em conteúdo de nicho, otimizando conjuntamente \textit{cache} e balanceamento.

\subsection{Análise de Custos Operacionais}

Estudo econômico comparando custos operacionais (energia, largura de banda, depreciação de \textit{hardware}) sob diferentes estratégias de balanceamento. Balanceamento ótimo pode não ser aquele que maximiza QoE absoluta, mas que maximiza QoE por unidade de custo, considerando realidades econômicas de operação de CDN.

\subsection{Content Steering com Múltiplas CDNs}

Extensão do algoritmo de monitoramento de memória para cenários de \textit{Content Steering} em MPEG-DASH, onde clientes podem alternar dinamicamente entre múltiplas CDNs durante uma sessão de \textit{streaming}. Este trabalho focou em balanceamento intra-CDN, operando dentro de um único Ponto de Presença (PoP). Uma direção promissora é aplicar princípios similares de monitoramento de recursos para decisões inter-CDN, onde cada CDN possui seu próprio balanceador de carga interno, mas um sistema de \textit{steering} de nível superior decide qual CDN utilizar com base em métricas agregadas de capacidade.

Content Steering, padronizado recentemente no DASH-IF, permite que clientes recebam instruções sobre quais CDNs priorizar através de documentos de steering atualizados dinamicamente. Integrar monitoramento de memória e carga neste contexto permitiria \textit{steering content-aware} e \textit{resource-aware}, direcionando tráfego não apenas baseado em proximidade geográfica ou custo, mas também em disponibilidade real de recursos em cada CDN. Isto é especialmente relevante para provedores de \textit{streaming} que utilizam múltiplas CDNs simultaneamente para redundância e otimização de custos.

\subsection{Direct Server Return (DSR)}

Investigação de arquiteturas de balanceamento baseadas em \textit{Direct Server Return}, onde o balanceador de carga processa apenas requisições de entrada (caminho de ida), mas respostas são enviadas diretamente dos servidores aos clientes, evitando o balanceador no caminho de retorno. Em \textit{streaming} de vídeo, onde respostas (segmentos de vídeo) são substancialmente maiores que requisições, DSR pode reduzir dramaticamente carga no balanceador e latência percebida.

Aplicar monitoramento de memória em arquitetura DSR introduz desafios interessantes: o balanceador não observa o tráfego de resposta, não podendo medir diretamente taxa de transferência ou sucesso de entregas. Seria necessário confiar exclusivamente em métricas do lado do servidor (memória, disco, \textit{CPU}) e potencialmente informações agregadas reportadas periodicamente. Avaliar trade-offs entre benefícios de DSR (menor latência, maior vazão) e complexidade de monitoramento é direção valiosa, especialmente para CDNs de ultra-larga escala onde o próprio balanceador pode tornar-se gargalo.

\subsection{\textit{Software}-Defined Networking (SDN)}

Exploração de como princípios de \textit{Software-Defined Networking} podem potencializar balanceamento dinâmico baseado em recursos. Em arquiteturas SDN, \textit{switches} são programáveis através de controladores centralizados, permitindo reconfiguração de fluxos de tráfego em tempo real sem intervenção manual. Este paradigma alinha-se naturalmente com balanceamento adaptativo baseado em monitoramento.

Uma direção promissora é instrumentar \textit{switches} SDN com informações de carga dos servidores, permitindo decisões de encaminhamento no próprio switch baseadas em métricas atuais. O controlador SDN poderia receber métricas de memória e disco dos servidores (via MQTT ou outro protocolo), calcular pesos probabilísticos usando o algoritmo proposto, e programar tabelas de fluxo nos \textit{switches} para distribuir tráfego proporcionalmente. Isto move decisões de balanceamento para mais próximo da borda da rede, potencialmente reduzindo latência de decisão e eliminando saltos adicionais através de balanceadores dedicados.

Adicionalmente, SDN permite experimentação com esquemas de balanceamento sofisticados como \textit{flow-based routing} (diferentes fluxos para diferentes servidores) versus \textit{packet-based routing} (cada pacote pode ir para servidor diferente, requerendo estados complexos), e \textit{elephant vs. mice flow handling} (grandes sessões de \textit{streaming} tratadas diferentemente de pequenas requisições). A integração de monitoramento de recursos com SDN representa fronteira empolgante para balanceamento de carga de próxima geração.

\section{Reprodutibilidade dos Experimentos}

Todo o código-fonte, scripts de configuração, infraestrutura experimental, e dados utilizados estão disponibilizados publicamente em repositório Git no endereço:

\begin{center}
\url{https://github.com/peedrofernandes/memory-oriented-load-balancer.git}
\end{center}

O repositório está organizado de forma a facilitar a reprodução exata de cada cenário experimental.

\begin{itemize}
    \item Servidor MPEG-DASH implementado em \textit{C\#} com \textit{ASP.NET Core}
    \item Cliente de validação com visualização em tempo real
    \item Gerador de carga parametrizável em \textit{Python}
    \item Sistema de monitoramento de recursos em \textit{Rust}
    \item Infraestrutura de comunicação via MQTT
    \item Arquivo de vídeo de teste codificado em múltiplas qualidades
    \item Scripts de orquestração com \textit{Docker Compose}
    \item Ferramentas de análise e geração de gráficos
\end{itemize}

Para reproduzir os experimentos específicos, foram criadas \textit{branches} dedicadas para cada combinação de cenário e estratégia de balanceamento, totalizando 16 configurações experimentais:

\begin{itemize}
    \item \textbf{Cenário 1} (100 usuários, servidores heterogêneos):
    \begin{itemize}
        \item \texttt{scenario-1-memory-monitoring-6s}
        \item \texttt{scenario-1-memory-monitoring-30s}
        \item \texttt{scenario-1-round-robin}
        \item \texttt{scenario-1-random-selection}
    \end{itemize}
    
    \item \textbf{Cenário 2} (100 usuários, servidores homogêneos):
    \begin{itemize}
        \item \texttt{scenario-2-memory-monitoring-6s}
        \item \texttt{scenario-2-memory-monitoring-30s}
        \item \texttt{scenario-2-round-robin}
        \item \texttt{scenario-2-random-selection}
    \end{itemize}
    
    \item \textbf{Cenário 3} (1000 usuários, servidores heterogêneos):
    \begin{itemize}
        \item \texttt{scenario-3-memory-monitoring-6s}
        \item \texttt{scenario-3-memory-monitoring-30s}
        \item \texttt{scenario-3-round-robin}
        \item \texttt{scenario-3-random-selection}
    \end{itemize}
    
    \item \textbf{Cenário 4} (1000 usuários, servidores homogêneos):
    \begin{itemize}
        \item \texttt{scenario-4-memory-monitoring-6s}
        \item \texttt{scenario-4-memory-monitoring-30s}
        \item \texttt{scenario-4-round-robin}
        \item \texttt{scenario-4-random-selection}
    \end{itemize}
\end{itemize}

Cada \textit{branch} contém configurações pré-ajustadas (arquivos \texttt{docker-compose.yml}, parâmetros de gerador de carga, configurações de servidores) correspondentes exatamente às condições experimentais descritas no Capítulo~\ref{cap:Resultados}. Para reproduzir um cenário específico, basta fazer \textit{checkout} da \textit{branch} correspondente e executar os scripts de orquestração fornecidos.

Esta organização permite não apenas validação independente dos resultados apresentados, mas também facilita experimentação com variações dos cenários (diferentes números de usuários, outras configurações de heterogeneidade, intervalos de atualização alternativos) ou adaptação do sistema para outros contextos de pesquisa. A documentação completa, incluindo instruções detalhadas de execução e requisitos de sistema, está disponível no arquivo \texttt{README.md} do repositório.

\section{Considerações Finais}

Este trabalho demonstrou de forma conclusiva que o monitoramento inteligente de recursos, especificamente memória principal e taxa de leitura de disco, possibilita melhorias substanciais em Qualidade de Experiência para \textit{streaming} de vídeo adaptativo em Redes de Distribuição de Conteúdo com servidores heterogêneos. As reduções de latência de até 85,7\%, eliminação de travamentos em cenários críticos, e manutenção de qualidade de vídeo alta e consistente validam a abordagem proposta e demonstram sua superioridade sobre estratégias convencionais em contextos apropriados.

A identificação da heterogeneidade como fator crítico determina claramente quando investir em balanceamento inteligente: essencial em CDNs reais com servidores de múltiplas gerações e capacidades, mas de valor marginal em ambientes homogêneos padronizados. A robustez demonstrada com dados moderadamente desatualizados resolve uma preocupação fundamental sobre escalabilidade, mostrando que o algoritmo pode operar eficientemente mesmo com intervalos de atualização maiores, reduzindo custos de monitoramento proporcionalmente.

Do ponto de vista prático, este trabalho oferece não apenas validação acadêmica de conceitos, mas uma solução implementável e testada que pode ser adaptada para sistemas de produção. A implementação de código aberto, infraestrutura experimental reutilizável, e análise detalhada de diferentes cenários fornecem base sólida para profissionais da indústria avaliarem aplicabilidade em seus contextos específicos.

Para a indústria de \textit{streaming} de vídeo, que movimenta bilhões de dólares globalmente e cujo crescimento é impulsionado por demanda crescente por conteúdo em alta qualidade, qualquer melhoria em QoE traduz-se diretamente em maior satisfação de usuários, redução de abandono (\textit{churn}), e vantagem competitiva. Os resultados deste trabalho sugerem que provedores de CDN operando em ambientes heterogêneos — a grande maioria — podem beneficiar significativamente de balanceamento baseado em monitoramento de recursos.

Olhando para o futuro, o \textit{streaming} de vídeo continua evoluindo: resolução \textit{4K}/\textit{8K}, realidade virtual, transmissão ao vivo interativa, e novos \textit{codecs} (\textit{AV1}, \textit{VVC}) aumentam complexidade e requisitos de recursos. Simultaneamente, \textit{edge computing} e redes \textit{5G} distribuem conteúdo ainda mais próximo aos usuários, criando ambientes inerentemente heterogêneos com milhares de pontos de presença de capacidades variadas. Neste contexto, balanceamento inteligente e adaptativo não é luxo opcional, mas necessidade fundamental.

Este trabalho estabelece fundação sólida, mas há muito a explorar. As direções de trabalho futuro delineadas — desde implantação em produção até aprendizado de máquina preditivo — representam apenas o início de uma jornada mais longa rumo a sistemas de distribuição de conteúdo verdadeiramente inteligentes, que aprendem, adaptam-se, e otimizam-se continuamente para proporcionar a melhor experiência possível aos usuários, independentemente de onde estejam ou quais recursos estejam disponíveis.

A mensagem central permanece clara: em um mundo onde conteúdo de vídeo domina o tráfego da Internet e expectativas de usuários por qualidade aumentam continuamente, decisões informadas por dados sobre onde e como distribuir esse conteúdo não são apenas técnicas — são estratégicas. O balanceamento de carga baseado em monitoramento de recursos é um passo importante nessa direção, demonstrando que sistemas que observam, interpretam e respondem ao estado real de seus componentes superam aqueles que operam cegamente, especialmente quando confrontados com a heterogeneidade e variabilidade inerentes à infraestrutura moderna de Internet.
